% !TeX root = ../main.tex
\section{Влияние амплитудно-фазового разбаланса квадратурных каналов}
Вернем значение гетеродина равное номеру варианта $ n_{g} = 9 $ и выставим тип \texttt{nco\_type = 'single'}. Пронаблюдаем несколько случаев: когда у нас идеальный гетородин без отклонений (цифровой, так как на реальных компонентах всегда присутствует разброс (amv =0 и pmv = 0)), только амплитудное несоответсвие (в 1 дБ), только фазовая ошибка (в $5^\circ$) и их совместное рассогласование в каналах (и амплитудная и фазовая).

\begin{figure}[H]
    \centering
    \includegraphics[width=0.8\textwidth]{4_spec_a0_p0.png}
    \caption{Без рассогласований}
\end{figure}

\begin{figure}[H]
    \centering
    \begin{minipage}[t]{0.48\textwidth}
        \vspace{0pt}
        \centering
        \includegraphics[width=\linewidth]{4_spec_a10_p0.png}
        \caption{Амплитудный рассогласование (1 дБ)}
        \label{fig:4_amv}
    \end{minipage}
    \begin{minipage}[t]{0.48\textwidth}
        \vspace{0pt}
        \centering
        \includegraphics[width=\linewidth]{4_spec_a0_p5.png}
        \caption{Фазовый рассогласование ($5^\circ$)}
        \label{fig:4_pmv}
    \end{minipage}
    \hfill
\end{figure}

\begin{figure}[H]
    \centering
    \includegraphics[width=0.8\textwidth]{4_spec_a15_p5.png}
    \caption{Рассогласование по амплитуде и фазе}
\end{figure}

Анализирую данные графики можно заметить, что в случае без рассогласований зеркальная составляющей не наблюдается, при амплитудном рассогласовании заметна ассиметрия относительно 0 Гц, при фазовом рассогласовании имеет место зеркальная составляющая, при обоих типах рассогласования мощность мешающей компоненты $P_r$ максимально из все приведённых опытов.